\documentclass[11pt]{article}

\title{Row Reduction and Canonical Matrix Forms}
\author{UlamLib seed note}
\date{}

\newtheorem{definition}{Definition}[section]
\newtheorem{lemma}[definition]{Lemma}
\newtheorem{theorem}[definition]{Theorem}
\newtheorem{corollary}[definition]{Corollary}

\begin{document}
\maketitle

\section{Scope}

This note is a seed target for a student-facing linear algebra layer over a
field. The intended Lean destination is
\texttt{ULAM/LinearAlgebra/RowReduction.lean}.

\section{Core definitions}

\begin{definition}
An elementary row operation on an $m \times n$ matrix over a field $K$ is one
of the following:
\begin{enumerate}
\item swap two rows,
\item multiply a row by a nonzero scalar,
\item add a scalar multiple of one row to another row.
\end{enumerate}
\end{definition}

\begin{definition}
A matrix is in row echelon form if every nonzero row lies above every zero row,
the leading nonzero entry of each nonzero row lies strictly to the right of the
leading entry of the previous nonzero row, and every entry below a leading
entry is zero.
\end{definition}

\begin{definition}
A matrix is in reduced row echelon form if it is in row echelon form, every
leading entry is equal to $1$, and every other entry in a pivot column is zero.
\end{definition}

\section{Seed theorem cluster}

\begin{lemma}
Each elementary row operation preserves row equivalence.
\end{lemma}

\begin{theorem}[Existence of echelon form]
Every matrix over a field is row-equivalent to a matrix in row echelon form.
\end{theorem}

\begin{theorem}[Existence of reduced row echelon form]
Every matrix over a field is row-equivalent to a matrix in reduced row echelon
form.
\end{theorem}

\begin{theorem}[Rank from pivots]
If a matrix is in reduced row echelon form, then its rank is equal to the
number of pivot columns.
\end{theorem}

\begin{corollary}
Every homogeneous linear system over a field admits a basis of solutions that
can be extracted from a reduced row echelon form computation.
\end{corollary}

\section{Formalization backlog}

\begin{enumerate}
\item Define a row and column operation DSL.
\item Define pivot positions and echelon predicates.
\item Implement Gaussian elimination with a certificate of row operations.
\item Package reduced row echelon form and rank computations.
\item Bridge the elimination API to canonical-form statements.
\end{enumerate}

\end{document}
